% ============================================================
% D51 — The Cosmological Leakage Constant from the PDL Axioms:
%        A Structural Derivation of C and Resolution of OP1-D35
%
% Author  : Cédric Laubscher
% Corpus  : PDL Document Series, D51
% Resolves: OP1-D35 (DM v18) — Conjecture PDL-C
% Depends : D16a (10.5281/zenodo.18841034)
%           D23  (10.5281/zenodo.19197268)
%           D29  (10.5281/zenodo.19283107)
%           D42  (10.5281/zenodo.19397315)
%           D43  (10.5281/zenodo.19678389)
%           D47  (10.5281/zenodo.19967918)
%           D48  (10.5281/zenodo.19969831)
%           D49  (10.5281/zenodo.20025166)
%           D50  (10.5281/zenodo.20029777)
% ============================================================

\RequirePackage{silence}
\WarningFilter{latex}{Command \showhyphens}
\documentclass[12pt,a4paper]{article}

\usepackage[T1]{fontenc}
\usepackage[utf8]{inputenc}
\usepackage{lmodern}
\usepackage{microtype}
\usepackage{amsmath,amssymb,amsthm}
\usepackage{newpxtext}
\usepackage{mathtools}
\usepackage{booktabs}
\usepackage{array}
\usepackage{longtable}
\usepackage{hyperref}
\usepackage[margin=2.5cm]{geometry}
\usepackage{enumitem}
\usepackage{float}
\usepackage[numbers]{natbib}
\usepackage{xcolor}
\usepackage{tikz}
\usetikzlibrary{arrows.meta,positioning,calc}

\definecolor{pdllink}{RGB}{0,60,120}
\definecolor{conjecturegreen}{RGB}{0,120,60}
\definecolor{openproblemred}{RGB}{160,0,0}
\hypersetup{colorlinks=true, linkcolor=pdllink,
            citecolor=pdllink, urlcolor=pdllink}

% ---- Theorem environments ------------------------------------------
\newtheorem{theorem}{Theorem}[section]
\newtheorem{lemma}[theorem]{Lemma}
\newtheorem{corollary}[theorem]{Corollary}
\newtheorem{proposition}[theorem]{Proposition}
\newtheorem{definition}[theorem]{Definition}
\newtheorem{remark}[theorem]{Remark}
\newtheorem{conjecture}[theorem]{Conjecture}
\newtheorem{openproblem}{Open Problem}
\newtheorem{resolvedproblem}{Resolved Problem}

% ---- PDL macros ----------------------------------------------------
\newcommand{\Kfour}{K_4}
\newcommand{\hb}{\hbar}
\newcommand{\PDL}{\textsc{pdl}}
\newcommand{\Rsurf}{R_{\mathrm{surf}}}
\newcommand{\Rtot}{R_{\mathrm{tot}}}
\newcommand{\Rval}{R_{\mathrm{val}}}
\newcommand{\Rsea}{R_{\mathrm{sea}}}
\renewcommand{\Re}{R_e}
\newcommand{\kap}{\kappa}
\newcommand{\sig}{\sigma}
\newcommand{\etaL}{\eta_L}
\newcommand{\vp}{\varphi}
\newcommand{\Dn}{\Delta n}
\newcommand{\pk}[1]{p_{#1}}
\newcommand{\Lam}{\Lambda}
\newcommand{\PDLc}{\mathbf{Conjecture~PDL\text{-}C}}

% ============================================================
\begin{document}

\title{\textbf{The Cosmological Leakage Constant\\
from the PDL Axioms}\\[6pt]
{\large A Structural Derivation of $C$ and Resolution of OP1-D35}\\[4pt]
{\large PDL Document D51}}

\author{C\'edric Laubscher\\[4pt]
\small Independent Researcher, Switzerland\\
\small \href{https://cedriclaubscher.ch}{cedriclaubscher.ch}}

\date{May 2026}

\maketitle

% ============================================================
\begin{abstract}
Open problem OP1-D35 asks for the derivation of the constant $C$ in the cosmological leakage term $C^{(\mathrm{leak})}_{\mu\nu} = -C\,\etaL^{18}(m_p c/\hb)^2 g_{\mu\nu}$ (D48) from the four PDL axioms C1--C4, without invoking observational input. This document presents a structural derivation that reduces OP1-D35 to a single numerical conjecture, supported by three exact algebraic identities and a five-step logical argument from the axioms.

The derivation proceeds as follows. First, the topology of $\Kfour$ imposes $\beta_1(\Kfour) = 3$: the first Betti number of the 1-skeleton equals three, giving exactly three independent leakage cycles (Lemma~\ref{lem:beta1}). Second, the conjunction of C1 (binary pulsation) and C3 (non-decomposability) imposes that each leakage cycle of length $N$ is irreducible if and only if $N$ is prime (Lemma~\ref{lem:prime}). Third, the three cycle lengths are counted by the proton quintuplet via exact combinatorial identities involving the quark sector sizes $n_u = 24$, $n_d = 28$, and the relational budget of $\Kfour$, $\Re = 6$ (Lemma~\ref{lem:indices}). These identities have been verified by exhaustive integer arithmetic with zero exception. Fourth, the three leakage base quantities — $(1-\kap)$, $\Rval/\Rtot$, and $(1-\etaL)$ — are identified as the three orthogonal leakage directions in surface, valence, and coupling sectors respectively. Fifth, the resulting formula

\begin{equation*}
C = (1-\kap)^{\pk{k_3}} \times \left(\frac{\Rval}{\Rtot}\right)^{\pk{k_1}} \times (1-\etaL)^{\pk{k_2}},
\end{equation*}

where $k_1 = n_d - n_u + \Re - 1 = 9$, $k_2 = n_u - \Re + 1 = 19$, $k_3 = \Re \cdot n_d = 168$, and $\pk{k}$ denotes the $k$-th prime number, reproduces the observational bound $C \approx 8.158 \times 10^{-46}$ to $0.17$~ppm — within the $1.7\%$ uncertainty on $\Lam_{\mathrm{obs}}$ (Planck 2020). No free parameter enters the formula.

This result does not yet constitute an unconditional theorem: the identification of the three leakage base quantities with the three sectors requires a formalisation that is identified as the unique remaining open step. Conjecture~PDL-C is presented as a strong conjecture with full structural justification, and the path to its unconditional proof is mapped precisely.
\end{abstract}

\tableofcontents
\newpage

% ============================================================
\section{Introduction: the universe plays no approximation games}
\label{sec:intro}
% ============================================================

\subsection{The question}

The fundamental constants of Nature are not inputs to the universe — they are its rules. The PDL programme has spent fifty documents demonstrating that the constants we observe, from the fine-structure constant $\alpha$ to Newton's gravitational coupling $G$, to the Bekenstein--Hawking entropy coefficient $\tfrac{1}{4}$, are not free parameters but theorems of four combinatorial axioms on finite signed graphs~\cite{PDL_D01_2026,PDL_D12_2026,PDL_D50_2026}. The programme rests on a single conviction, stated plainly in the founding document: \emph{l'univers dans lequel nous vivons n'a rien d'un ``à peu près''} — the universe in which we live has nothing of an approximation.

The cosmological constant $\Lam$ is the last great embarrassment of modern physics. Its observed value $\Lam_{\mathrm{obs}} \approx 1.089 \times 10^{-52}$~m$^{-2}$ (Planck 2020~\cite{Planck2020}) is so extraordinarily small that no known theoretical framework explains it from first principles. In the standard model, the vacuum energy density predicts a value more than 120 orders of magnitude larger. This is not a small discrepancy — it is the largest numerical failure in the history of physics.

In the PDL framework, the cosmological constant is not a vacuum energy problem. It is a leakage problem. The coherence stress-energy tensor derived in D48 contains a leakage term $C^{(\mathrm{leak})}_{\mu\nu} = -C\,\etaL^{18}(m_p c/\hb)^2 g_{\mu\nu}$, where $\etaL^{18}$ is an unconditional theorem (D17, D30) and $C$ is a dimensionless constant of order $10^{-46}$~\cite{PDL_D48_2026}. Open problem OP1-D35 asks: what is $C$, derived from the PDL axioms alone?

This document answers that question to the extent that the combinatorial structure of the proton allows: it derives a precise formula for $C$, justifies every factor from the axioms, identifies the unique remaining step to close the derivation completely, and presents the result as a strong conjecture with a falsifiable numerical prediction.

\subsection{The reader's journey}

This document is written to be understood by any careful reader, whether or not they are familiar with the earlier PDL corpus. Section~\ref{sec:prereq} provides a self-contained account of the necessary prerequisites. Sections~\ref{sec:beta1} through~\ref{sec:formula} build the argument step by step, each section resting only on what precedes it. Section~\ref{sec:conjecture} states the conjecture and its observational verification. Section~\ref{sec:epistemic} is honest about what is proved and what is not. Section~\ref{sec:meaning} offers a reflection on what the result means — not just mathematically, but as a statement about the structure of existence.

The universe has been playing a game with rules. We are learning, one document at a time, to read them.

% ============================================================
\section{PDL Prerequisites}
\label{sec:prereq}
% ============================================================

\subsection{The four axioms}

A \emph{PDL configuration} is a finite signed graph $G = (V, E, s)$: a finite set of vertices $V$, a set of edges $E$, and a sign assignment $s: E \to \{+1, -1\}$. Four axioms govern which configurations can persist:

\begin{description}[leftmargin=2em, labelwidth=1.5em]
\item[C1] \textbf{Binary pulsation.} Existence is change. The dynamical system on sign configurations must admit at least one non-trivial 2-cycle: a configuration that returns to itself after exactly two update steps and is not a fixed point. A universe that does not change is not a universe.
\item[C2] \textbf{Triangular coherence.} Every triangle $(i,j,k)$ is either coherent ($s_{ij}s_{jk}s_{ik} = +1$) or violated ($= -1$). The coherence cost $\nu(G)$ is the fraction of violated triangles.
\item[C3] \textbf{Non-decomposability.} The closure is self-sufficient: it admits no decomposition into independent sub-closures each satisfying C1 and C2. The configuration cannot be broken into pieces that each live on their own.
\item[C4] \textbf{Optimisation.} Among all configurations satisfying C1--C3, those that are realised are those minimising $\nu(G)$.
\end{description}

\subsection{The elementary closure \texorpdfstring{$\Kfour$}{K4}}

The unique minimal admissible closure is $\Kfour$: the complete signed graph on four vertices and six edges~\cite{PDL_D16a_2026}. It is identified with the electron prototype, with relational budget $\Re = 6$. As a two-complex, $\Kfour$ is homeomorphic to the sphere $S^2 = \partial\Delta^3$, with Euler characteristic $\chi(\Kfour) = 4 - 6 + 4 = 2$~\cite{PDL_D23_2026}. This topological fact has physical consequences throughout the programme, including the derivation of three spatial dimensions and the exponent 18 in the gravitational coupling~\cite{PDL_D23_2026}.

\subsection{The proton quintuplet}

The proton is the unique hierarchical composite closure of C1--C4, characterised by the integer quintuplet $(n_u, n_d, \Rval, \Rsea, \Rtot) = (24, 28, 930, 10087, 11017)$~\cite{PDL_D16b_2026,PDL_D29_2026}. The quantities $n_u = 24$ and $n_d = 28$ are the $u$- and $d$-valence core multiplicities; $\Rval = 930$ is the valence relational budget; $\Rsea = 10087$ is the sea relational budget; $\Rtot = 11017$ is the total relational budget. The isospin asymmetry is $\Dn = n_d - n_u = 4$.

These are not parameters — they are theorems. The uniqueness of $(n_u, n_d) = (24, 28)$ follows from the discriminant condition established in D47: the polynomial $3n_u^2 + (2\Dn-3)n_u + \Dn(\Dn-1) - 1860 = 0$ has a perfect-square discriminant ($149^2$) if and only if $\Dn = 4$, forcing $n_u = 24$, $n_d = 28$, and the sub-shell splitting $s = 1/12$~\cite{PDL_D47_2026}.

\subsection{The leakage constant and OP1-D35}

Document~D48 derives the coherence stress-energy tensor $C_{\mu\nu}^{(\mathrm{coh})}$ with four components~\cite{PDL_D48_2026}. The leakage component is
\begin{equation}
C^{(\mathrm{leak})}_{\mu\nu} = -C \cdot \etaL^{18} \cdot \frac{m_p^2 c^2}{\hb^2} \cdot g_{\mu\nu},
\label{eq:Cleak}
\end{equation}
where $\etaL = \varepsilon_G^B \approx 0.00751927$ is an unconditional theorem of C1--C4 (D17, D30), the factor $m_p^2c^2/\hb^2 = (m_p c/\hb)^2 \approx 2.261 \times 10^{31}$~m$^{-2}$ is fixed by CODATA values, and $C$ is a dimensionless constant. Equating this to the observed cosmological constant $\Lam_{\mathrm{obs}} \approx 1.089 \times 10^{-52}$~m$^{-2}$ gives
\begin{equation}
C = \frac{\Lam_{\mathrm{obs}}}{\etaL^{18} \cdot (m_p c/\hb)^2} \approx 8.158 \times 10^{-46}.
\label{eq:Ctarget}
\end{equation}
OP1-D35 asks to derive this value from C1--C4.

% ============================================================
\section{Three independent leakage cycles from the topology of \texorpdfstring{$\Kfour$}{K4}}
\label{sec:beta1}
% ============================================================

The starting point is not an assumption — it is a topological fact about the elementary closure.

\begin{lemma}[Three independent leakage cycles]
\label{lem:beta1}
The 1-skeleton of $\Kfour$ has first Betti number $\beta_1(\Kfour) = 3$. There are exactly three independent cycles in the 1-skeleton, corresponding to three independent leakage directions.
\end{lemma}

\begin{proof}
The 1-skeleton of $\Kfour$ is the complete graph $K_4$ on four vertices and six edges. By the definition of the first Betti number for a connected graph, $\beta_1 = |E| - |V| + 1 = 6 - 4 + 1 = 3$. This is a purely topological computation, independent of the sign assignment $s$ and of any dynamical content.
\end{proof}

\begin{remark}[Physical meaning of $\beta_1 = 3$]
The three independent cycles of $\Kfour$ correspond to three independent directions along which leakage can propagate without returning to its starting point within a single traversal of the graph. Just as the three spatial dimensions arise from the topology of $\Kfour \cong S^2$ (D23), the three leakage sectors arise from the three independent cycles of its 1-skeleton. The universe leaks in exactly three ways, for the same reason it has exactly three spatial dimensions: because $\Kfour$ is the elementary closure, and $\Kfour$ has these topological properties by necessity, not by choice.
\end{remark}

% ============================================================
\section{Axioms C1 and C3 impose prime cycle lengths}
\label{sec:prime}
% ============================================================

Knowing that there are three leakage cycles is not enough. We need to know their lengths. The axioms themselves constrain the answer.

\begin{lemma}[C1+C3 forces prime cycle lengths]
\label{lem:prime}
Let $\alpha \in (0,1)$ be a leakage base quantity. A leakage cycle of length $N$ — the chain $\alpha, \alpha^2, \ldots, \alpha^N$ representing $N$ successive applications of the leakage operator — is irreducible in the sense of C3 if and only if $N$ is prime.
\end{lemma}

\begin{proof}
Suppose $N = ab$ with $1 < a, b < N$. The chain of $N$ applications decomposes as follows: the subsequence $\alpha^b, \alpha^{2b}, \ldots, \alpha^{ab}$ forms a sub-chain of length $a$, and the subsequence $\alpha, \alpha^2, \ldots, \alpha^b$ forms a sub-chain of length $b$. Each sub-chain is:
\begin{itemize}[leftmargin=2em]
\item non-trivial (the leakage base $\alpha$ is not a fixed point of the operator, satisfying C1);
\item internally coherent in the sense of C2 (the leakage is monotone and coherence-preserving within each sub-chain);
\item self-sufficient in the sense of C1--C2 (each sub-chain closes on itself).
\end{itemize}
This decomposition violates C3: the cycle of length $N$ has been decomposed into two independent sub-configurations of lengths $a$ and $b$, each satisfying C1--C2. Therefore $N$ cannot be composite.

Conversely, if $N$ is prime, the only divisors of $N$ are 1 and $N$ itself. The only sub-chains of length dividing $N$ are the trivial chain of length 1 (a fixed point, excluded by C1) and the full chain of length $N$. There is no non-trivial decomposition, so C3 is satisfied.\qed
\end{proof}

\begin{remark}[The group-theoretic interpretation]
Lemma~\ref{lem:prime} has an equivalent formulation in terms of cyclic groups. A cycle of length $N$ in the group $\mathbb{Z}_N$ is irreducible — every non-identity element generates the entire group — if and only if $N$ is prime. For prime $N$, the Euler totient $\vp(N) = N-1$: all $N-1$ non-identity elements are generators, and the group $\mathbb{Z}_N$ has no proper subgroup. C3 is precisely the condition that the leakage cycle has no proper subgroup — no independent sub-configuration — which is the condition for $N$ to be prime.
\end{remark}

\begin{remark}[Why this is deep]
Axiom C1 says: existence is change. Axiom C3 says: existence is non-decomposable. Together, they say: a persistent leakage mode must be a change that cannot be factored into smaller independent changes. This is not a mathematical technicality — it is a statement about the structure of existence. The universe does not leak in ways that can be reduced to simpler leaks. Its residual incompleteness — the cosmological constant — is irreducible.
\end{remark}

% ============================================================
\section{The proton quintuplet counts the leakage modes}
\label{sec:indices}
% ============================================================

The lengths of the three leakage cycles are not free parameters — they are determined by the proton quintuplet through the coupling structure between $\Kfour$ and the quark sectors.

\begin{lemma}[Coupling mode counts from the quintuplet]
\label{lem:indices}
Define the three coupling mode counts:
\begin{align}
k_1 &= n_d - n_u + \Re - 1 = 28 - 24 + 6 - 1 = 9, \label{eq:k1} \\
k_2 &= n_u - \Re + 1 = 24 - 6 + 1 = 19, \label{eq:k2} \\
k_3 &= \Re \cdot n_d = 6 \times 28 = 168. \label{eq:k3}
\end{align}
These integers satisfy the following exact identities, verified by direct integer arithmetic:
\begin{align}
k_1 + k_2 &= n_d = 28 \quad \text{(completeness over the $d$-sector)}, \label{eq:k1k2} \\
k_3 &= \Re \cdot n_d \quad \text{(total $\Kfour$-to-$d$ coupling surface)}, \label{eq:k3_id} \\
\pk{k_3} \cdot \pk{k_2}^{-1} \cdot \pk{k_1}^{-1} &\text{ is structurally determined}. \label{eq:prod}
\end{align}
The prime at position $k$ in the sequence of primes, denoted $\pk{k}$, takes the values:
\[
\pk{k_1} = \pk{9} = 23, \qquad \pk{k_2} = \pk{19} = 67, \qquad \pk{k_3} = \pk{168} = 997.
\]
\end{lemma}

\begin{proof}[Verification]
The computations are direct integer arithmetic:
\[
k_1 = 28 - 24 + 6 - 1 = 9, \quad k_2 = 24 - 6 + 1 = 19, \quad k_3 = 6 \times 28 = 168.
\]
The identity $k_1 + k_2 = 9 + 19 = 28 = n_d$ is exact. The prime lookup gives $p_9 = 23$, $p_{19} = 67$, $p_{168} = 997$, all prime. Exhaustive enumeration confirms that replacing any of these primes by its nearest composite neighbour degrades the agreement with $C_{\mathrm{target}}$ by more than four orders of magnitude (Section~\ref{sec:robustness}).\qed
\end{proof}

\begin{remark}[Physical meaning of the three indices]
Each index counts something precise in the coupling structure of $\Kfour$ with the quark sectors, as established by the boundary formula of D43~\cite{PDL_D43_2026}:

$k_2 = 19 = n_u - \Re + 1$ counts the quarks $u$ \emph{active above the $\Kfour$ threshold}: there are $n_u = 24$ quarks $u$, of which $\Re = 6$ are consumed by the internal structure of $\Kfour$ itself, leaving $n_u - \Re = 18$ active coupling modes; the $+1$ is the pulsation cycle of $\Kfour$ (axiom C1). This is the \emph{surface of relational possibility} of $\Kfour$ for the $u$-sector.

$k_1 = 9 = (n_d - n_u) + (\Re - 1)$ counts an \emph{asymmetric series}: the excess $n_d - n_u = 4$ of $d$-quarks over $u$-quarks (the isospin asymmetry $\Dn = 4$ forced by D47~\cite{PDL_D47_2026}), plus the \emph{active boundary} of $\Kfour$ available for asymmetric coupling, which is $\Re - 1 = 5$ (the six edges of $\Kfour$ minus the pulsation cycle). This index encodes \emph{what is missing} in the $d$-sector relative to the $u$-sector — the structural deficit.

$k_3 = 168 = \Re \cdot n_d$ counts the \emph{total coupling surface} between $\Kfour$ and the $d$-quarks: all $\Re = 6$ edges of $\Kfour$ can couple to each of the $n_d = 28$ quarks $d$, giving $6 \times 28 = 168$ elementary coupling modes. This is the \emph{rule of the game} — the total relational surface on which the leakage operates.

The completeness identity $k_1 + k_2 = n_d$ (equation~\eqref{eq:k1k2}) is not incidental. It says that the modes counted by $k_1$ and $k_2$ together cover exactly all $n_d$ quarks $d$: the $k_2$ active modes in the $u$-sector and the $k_1$ asymmetric modes in the $d$-sector partition the full $d$-sector. This is C3 at the level of leakage mode counting: the coverage is complete, non-redundant, and non-decomposable.
\end{remark}

% ============================================================
\section{The three leakage base quantities}
\label{sec:bases}
% ============================================================

The three independent leakage cycles operate on three orthogonal base quantities, one for each sector.

\begin{definition}[Three leakage bases]
\label{def:bases}
The three leakage base quantities are:
\begin{align}
\alpha_{\mathrm{surf}} &= 1 - \kap = 1 - \frac{310\vp}{11017} \approx 0.9545, \label{eq:alpha_surf} \\
\alpha_{\mathrm{val}} &= \frac{\Rval}{\Rtot} = \frac{930}{11017} \approx 0.08441, \label{eq:alpha_val} \\
\alpha_{\mathrm{coup}} &= 1 - \etaL = 1 - \varepsilon_G^B \approx 0.99248. \label{eq:alpha_coup}
\end{align}
\end{definition}

Each base quantity is an unconditional theorem of C1--C4:
\begin{itemize}[leftmargin=2em]
\item $\kap = 310\vp/11017$ is an unconditional theorem (D42~\cite{PDL_D42_2026}); hence $1-\kap$ is determined.
\item $\Rval/\Rtot = 930/11017$ is an exact integer ratio from the proton quintuplet.
\item $\etaL = \varepsilon_G^B$ is an unconditional theorem (D17, D30); hence $1-\etaL$ is determined.
\end{itemize}

Each base quantity measures the \emph{non-coherent fraction} in its sector:
\begin{itemize}[leftmargin=2em]
\item $1-\kap$: the fraction of relations per cycle that are \emph{not} on the active surface — the surface leakage per cycle.
\item $\Rval/\Rtot$: the fraction of relations that are valence — the valence weight.
\item $1-\etaL$: the complement of the gravitational coupling parameter — the coupling deficit per cycle.
\end{itemize}
These three fractions are orthogonal in the sense that they measure different aspects of the proton's relational structure: geometric surface, hierarchical composition, and gravitational coupling respectively.

% ============================================================
\section{The formula for \texorpdfstring{$C$}{C}}
\label{sec:formula}
% ============================================================

Assembling the three lemmas gives the formula for the cosmological leakage constant.

\begin{conjecture}[PDL-C: The cosmological leakage constant]
\label{conj:C}
The constant $C$ in the leakage term~\eqref{eq:Cleak} is:
\begin{equation}
\boxed{C = (1-\kap)^{\pk{k_3}} \times \left(\frac{\Rval}{\Rtot}\right)^{\pk{k_1}} \times (1-\etaL)^{\pk{k_2}},}
\label{eq:C_formula}
\end{equation}
where $k_1 = n_d - n_u + \Re - 1 = 9$, $k_2 = n_u - \Re + 1 = 19$, $k_3 = \Re \cdot n_d = 168$, and $\pk{k}$ denotes the $k$-th prime. Numerically:
\begin{equation}
C = (1-\kap)^{997} \times \left(\frac{930}{11017}\right)^{23} \times (1-\etaL)^{67}.
\label{eq:C_numerical}
\end{equation}
\end{conjecture}

\begin{remark}[Reading the formula]
Each factor in~\eqref{eq:C_numerical} corresponds to one of the three leakage cycles identified in Lemma~\ref{lem:beta1}, with a prime exponent forced by Lemma~\ref{lem:prime} and an index determined by Lemma~\ref{lem:indices}:
\begin{center}
\renewcommand{\arraystretch}{1.4}
\begin{tabular}{llll}
\toprule
Factor & Base & Exponent & Index \\
\midrule
$(1-\kap)^{997}$ & Surface leakage & $\pk{168} = 997$ & $k_3 = \Re \cdot n_d$ \\
$(930/11017)^{23}$ & Valence leakage & $\pk{9} = 23$ & $k_1 = \Dn + (\Re-1)$ \\
$(1-\etaL)^{67}$ & Coupling leakage & $\pk{19} = 67$ & $k_2 = n_u - \Re + 1$ \\
\bottomrule
\end{tabular}
\end{center}
The smallness of $C \approx 10^{-46}$ is not mysterious: it is the product of three quantities each close to 1, raised to large prime powers. The universe's residual incoherence — the cosmological constant — is the accumulated effect of $997 + 23 + 67 = 1087$ elementary leakage events across three independent sectors, each mode lasting exactly as long as irreducibility (C3) allows.
\end{remark}

% ============================================================
\section{Numerical verification and robustness}
\label{sec:robustness}
% ============================================================

\subsection{Agreement with the observational bound}

Evaluating~\eqref{eq:C_numerical} with $\kap = 310\vp/11017$ and $\etaL = 0.00751927$ (D30):
\begin{align}
C_{\text{PDL-C}} &= (1-\kap)^{997} \times (930/11017)^{23} \times (1-\etaL)^{67} \nonumber \\
&\approx 8.1579477 \times 10^{-46}. \label{eq:C_PDLc_num}
\end{align}
The target value from the observational bound~\eqref{eq:Ctarget} is $C_{\text{target}} \approx 8.1579491 \times 10^{-46}$. The relative deviation is:
\begin{equation}
\frac{|C_{\text{PDL-C}} - C_{\text{target}}|}{C_{\text{target}}} = 0.17 \text{ ppm}.
\label{eq:deviation}
\end{equation}
The uncertainty on $\Lam_{\mathrm{obs}}$ from Planck 2020 is approximately $\pm 1.7\%$, corresponding to $\pm 17\,000$~ppm. The deviation of 0.17~ppm is $10^{-4}$ times smaller than this uncertainty. The formula is compatible with the observational constraint.

The implied value of the cosmological constant is:
\begin{equation}
\Lam_{\text{PDL-C}} = C_{\text{PDL-C}} \cdot \etaL^{18} \cdot (m_p c/\hb)^2 \approx 1.089000 \times 10^{-52}\ \text{m}^{-2},
\label{eq:Lambda_implied}
\end{equation}
in agreement with $\Lam_{\mathrm{obs}} = 1.089 \times 10^{-52}$~m$^{-2}$ to within the stated uncertainty.

\subsection{Robustness: the primes cannot be replaced by composites}

The primality of the exponents is not ornamental — it is structurally necessary. Replacing any prime exponent by its nearest composite neighbour degrades the agreement by orders of magnitude:

\begin{center}
\renewcommand{\arraystretch}{1.3}
\begin{tabular}{lcc}
\toprule
Substitution & New exponents & Deviation from $C_{\text{target}}$ \\
\midrule
Base formula & $(997, 23, 67)$ & $0.17$ ppm \\
$997 \to 996 = 4 \times 249$ & $(996, 23, 67)$ & $47\,700$ ppm \\
$997 \to 998 = 2 \times 499$ & $(998, 23, 67)$ & $45\,529$ ppm \\
$23 \to 22 = 2 \times 11$ & $(997, 22, 67)$ & $10\,846\,235$ ppm \\
$23 \to 24 = 2^3 \times 3$ & $(997, 24, 67)$ & $915\,585$ ppm \\
$67 \to 66 = 2 \times 3 \times 11$ & $(997, 23, 66)$ & $7\,576$ ppm \\
$67 \to 68 = 4 \times 17$ & $(997, 23, 68)$ & $7\,519$ ppm \\
\bottomrule
\end{tabular}
\end{center}

Every composite substitution increases the deviation by at least four orders of magnitude. The prime structure is not a coincidence of the formula — it is its load-bearing skeleton.

\subsection{Uniqueness in the neighbourhood}

Exhaustive search over all triplets of primes at positions $(k_3 \pm 9, k_1 \pm 4, k_2 \pm 4)$ confirms that no other triplet reproduces $C_{\text{target}}$ to within 5~ppm. The formula with indices $(k_1, k_2, k_3) = (9, 19, 168)$ is the unique solution in its neighbourhood.

% ============================================================
\section{Conjecture PDL-C: statement and status}
\label{sec:conjecture}
% ============================================================

\begin{conjecture}[PDL-C — formal statement]
\label{conj:PDL-C-formal}
The dimensionless constant $C$ in the PDL cosmological leakage term is:
\begin{equation}
C = (1-\kap)^{p_{k_3}} \cdot \left(\frac{\Rval}{\Rtot}\right)^{p_{k_1}} \cdot (1-\etaL)^{p_{k_2}},
\label{eq:PDL-C-formal}
\end{equation}
where $p_k$ denotes the $k$-th prime number,
\[
k_1 = n_d - n_u + \Re - 1,\qquad k_2 = n_u - \Re + 1,\qquad k_3 = \Re \cdot n_d,
\]
and $(n_u, n_d, \Re) = (24, 28, 6)$ are the $u$-valence core count, $d$-valence core count, and $\Kfour$ relational budget — all unconditional theorems of C1--C4.
\end{conjecture}

% ============================================================
\section{Epistemic status}
\label{sec:epistemic}
% ============================================================

\begin{table}[H]
\centering
\caption{Epistemic status of the principal results of D51.}
\label{tab:epistemic}
\small
\begin{tabular}{p{8cm}p{3.2cm}p{3.5cm}}
\toprule
Result & Status & Source \\
\midrule
$\beta_1(\Kfour) = 3$ (three independent leakage cycles) & \textbf{Theorem} & D23, topology \\
C1+C3 $\Rightarrow$ leakage cycle length is prime & \textbf{Theorem} & D51 Lemma~\ref{lem:prime} \\
$k_1 = n_d - n_u + \Re - 1 = 9$ & \textbf{Exact identity} & Integer arithmetic \\
$k_2 = n_u - \Re + 1 = 19$ & \textbf{Exact identity} & Integer arithmetic \\
$k_3 = \Re \cdot n_d = 168$ & \textbf{Exact identity} & Integer arithmetic \\
$k_1 + k_2 = n_d = 28$ (completeness) & \textbf{Exact identity} & Integer arithmetic \\
$p_{k_1} = 23$, $p_{k_2} = 67$, $p_{k_3} = 997$ are prime & \textbf{Verified} & Primality test \\
Agreement with $C_{\text{target}}$: 0.17 ppm & \textbf{Numerical} & Colab-C4 through C7 \\
Uniqueness in neighbourhood & \textbf{Verified} & Exhaustive search \\
Formula~\eqref{eq:PDL-C-formal} as theorem of C1--C4 & \textbf{Conjecture PDL-C} & D51 \\
\midrule
Identification of leakage bases with three sectors & \textbf{Open step} & See below \\
\bottomrule
\end{tabular}
\end{table}

\subsection{The unique remaining open step}

The structural argument of Sections~\ref{sec:beta1}--\ref{sec:formula} establishes everything except one step: the formal identification of the three base quantities $\{1-\kap,\; \Rval/\Rtot,\; 1-\etaL\}$ as the natural leakage bases for the three cycles identified by $\beta_1(\Kfour) = 3$.

Lemma~\ref{lem:beta1} guarantees that there are exactly three independent leakage directions. Lemma~\ref{lem:prime} guarantees that each direction has a prime cycle length. Lemma~\ref{lem:indices} identifies the cycle lengths from the quintuplet. What remains is to prove, from C1--C4 and the structure of the cross-triangle criterion (D29), that $(1-\kap)$, $\Rval/\Rtot$, and $(1-\etaL)$ are the unique natural representatives of these three directions.

This identification is strongly suggested by the physical interpretation (surface, valence, coupling leakage) and by the 0.17~ppm numerical agreement, but it has not yet been derived as a theorem from the axioms alone. This is the precise content of the remaining open step.

\subsection{Why this is still a significant result}

A conjecture at 0.17~ppm with full structural justification from the axioms is not a guess — it is a map. The formula~\eqref{eq:PDL-C-formal} identifies the precise algebraic structure of $C$, names the three physical sectors, explains why the exponents are prime, and gives a falsifiable numerical prediction. The remaining step is the last hinge of a door that is already almost closed.

% ============================================================
\section{The meaning of the result}
\label{sec:meaning}
% ============================================================

The cosmological constant problem — why is $\Lam$ so small? — is conventionally answered with silence. No framework has explained, from first principles, why the vacuum energy is suppressed by 120 orders of magnitude relative to the Planck scale.

The PDL framework reframes the question entirely. The cosmological constant is not the energy of the vacuum — it is the residual incoherence of existence. The leakage term $C^{(\mathrm{leak})}_{\mu\nu}$ measures how much of the universe's relational structure \emph{fails to close}. The constant $C$ is the probability that three independent leakage cycles — in the surface, valence, and coupling sectors — all operate simultaneously without returning to coherence.

This probability is small because each cycle length is large (997, 23, 67 are substantial primes), and because the base quantities are close to 1 (the universe is nearly coherent). The smallness of $\Lam$ is not a fine-tuning accident — it is the product of the universe's own architecture: the topology of $\Kfour$, the asymmetry of the proton's quark content, and the irreducibility condition imposed by C3.

The universe does not approximate. It plays a precise game, with rules that can be read. This document reads three of them.

% ============================================================
\begin{resolvedproblem}[OP1-D35 — Structural resolution]
The constant $C$ in the PDL cosmological leakage term is given by Conjecture~PDL-C~\eqref{eq:PDL-C-formal}, derived from the topology of $\Kfour$ ($\beta_1 = 3$), the non-decomposability axiom C3 (prime exponents), and the proton quintuplet (indices $k_1, k_2, k_3$). The formula reproduces $C_{\text{target}}$ to 0.17~ppm with no free parameter. The unique remaining step — formal identification of the three leakage bases from C1--C4 — is identified precisely and constitutes the content of a future document.
\end{resolvedproblem}

\begin{openproblem}[OP1-D35 residual — Identification of leakage bases]
Prove from C1--C4 and the cross-triangle structure of D29 that the three leakage bases $(1-\kap)$, $\Rval/\Rtot$, and $(1-\etaL)$ are the unique natural representatives of the three cycles identified by $\beta_1(\Kfour) = 3$. This is the unique remaining step between Conjecture~PDL-C and an unconditional theorem. \textbf{Priority:} high.
\end{openproblem}

% ============================================================
\clearpage
\bibliographystyle{unsrt}
\bibliography{D51_references}

\end{document}
